\documentclass[11pt]{article}
\usepackage{amsmath,amssymb,amsthm}
\newtheorem{proposition}{Proposition}
\begin{document}
\title{Restrictions on nesting second-order AC subspaces inside the explicit first-order family}
\author{}\date{}\maketitle

Put $j=N/2$ and $c=j(j+1)/3$. Let $V$ be the isometry whose columns
are the explicit maximal first-order AC basis in the catalogue, and let
$B$ be an isometry parametrizing a child subspace: $U=VB$. Since
$V^\dagger J_\alpha V=0$, every such child is automatically first-order AC.
It is second-order AC precisely when every traceless quadrupole has zero
compression to it. In particular, for any real linear combination $Q$ of
traceless quadrupoles,
\[
 B^\dagger V^\dagger QVB=0.
\]
If $V^\dagger QV$ has inertia $(n_+,n_-,n_0)$, the child dimension obeys
\[
 \dim U\le \min(n_+,n_-)+n_0.
\tag{1}
\]
Indeed, quotienting the zero eigenspace leaves a totally isotropic
subspace of a nondegenerate Hermitian form; projection to either sign
subspace is injective. This is a restriction on the chosen parent,
not on arbitrary second-order AC subspaces.

\begin{proposition}
In the catalogue's explicit maximal first-order parent at $N=18$, every
second-order AC child has dimension at most two. At $N=24$, every such
child has dimension at most three. These bounds also apply to collective
rotations of the chosen parents.
\end{proposition}
\begin{proof}
Use the reflection modules
\[
 v_x=\sqrt{\frac{x+1}{2x+1}}|x\rangle+
     \sqrt{\frac{x}{2x+1}}|-x-1\rangle,
 \quad
 w_x=\sqrt{\frac{x}{2x+1}}|x+1\rangle-
     \sqrt{\frac{x+1}{2x+1}}|-x\rangle.
\]
In the ordered parent basis $|0\rangle,v_2,w_2,v_5,w_5,\ldots$, take
\[
 Q_{18}=J_z^2-cI+(J_x^2-J_y^2)+\{J_x,J_z\},
 \qquad
 Q_{24}=J_z^2-cI+2(J_x^2-J_y^2)+\{J_x,J_z\}.
\]
Each has zero compression on every second-order AC subspace.
The exact $LDL^T$ pivots of $V^\dagger Q_{18}V$ belong, respectively, to
\[
 (-31,-29),\ (15,16),\ (-4,-3),\ (91,92),\ (54,55),\ (40,41),\ (5,6).
\]
Those of $V^\dagger Q_{24}V$ belong to
\[
 (-53,-51),\ (231,232),\ (-34,-33),\ (107,108),\ (52,53),
 \ (-33,-32),\ (377,378),\ (78,79),\ (26,27).
\]
These enclosures are certified by the accompanying verifier using exact
rational interval arithmetic with outward rounding at every operation.
It reconstructs the physical spin matrices and the exact radical parent
bases before eliminating. No numerical eigenvalue signs are assumed.
Sylvester's law gives inertia $(5,2,0)$ and $(6,3,0)$, respectively.
Equation (1) proves the claims. Collective rotations preserve both AC
order and the existence of a child subspace.
\end{proof}

\begin{proposition}
For every $N\ge22$, no second-order AC subspace inside the catalogue's
explicit maximal first-order family attains the general upper bound
\[
 \min\left\{\left\lfloor\frac{N-2}{4}\right\rfloor,
             \left\lfloor\frac{N+5}{5}\right\rfloor\right\}.
\]
In fact, for $N\ge15$ every such child has dimension at most
\[
 2\left\lfloor\frac{j+2-r}{3}\right\rfloor,
 \qquad r=\frac{\sqrt{1+4j(j+1)/3}-1}{2}.
 \tag{2}
\]
Its asymptotic coefficient is $(1-1/\sqrt3)/3\simeq0.140883$.
\end{proposition}
\begin{proof}
Use $Q=J_z^2-cI$. On each reflection module its compression is
\[
 (x(x+1)-c)I_2.
\]
The module parameters are spaced by three and satisfy $x\le j-1$.
All remaining central blocks are supported on $|m|\le9/2$; for $N\ge15$
we have $c>81/4$, so their compressed $Q$ is strictly negative.
There are no cross terms between distinct blocks since $Q$ is diagonal
in the magnetic basis and their supports are disjoint.
The number of nonnegative module directions is therefore twice the
number of parameters $x\ge r$, at most
$2\lfloor(j+2-r)/3\rfloor$. Applying (1), even without using the negative
inertia count, proves (2).

Since $r>j/\sqrt3-1/2$, its right-hand side is strictly smaller than
\[
 \frac{1-1/\sqrt3}{3}N+\frac53.
\]
For $N\ge29$ this is smaller than $N/5$: at $N=29$, the required
inequality is $145>83\sqrt3$, certified by
$145^2-3\cdot83^2=358>0$; the difference then increases with $N$.
For $N\ge30$, the general upper bound is
$\lfloor(N+5)/5\rfloor=\lfloor N/5\rfloor+1>N/5$.
At $N=29$ both terms give six. Finally, direct exact evaluation of (2)
gives four at each $N=22,\ldots,28$, whereas the general bound is
$5,5,5,5,6,6,6$, respectively.
\end{proof}

These exclusions do not characterize all first-order AC subspaces.
In particular, the known unrestricted optima $K(18,2)=4$ and $K(24,2)=5$
cannot be contained in these specific maximal first-order parents.
Searching over different first-order parents remains a valid construction
strategy. No extension theorem asserting containment in this explicit
family is used or established here.
\end{document}
